\section{Introduction}
\label{sec:intro}
Average symptom improvement does not settle whether cognitive behavioral therapy (CBT) adequately serves the people offered it.
A systematic review of anxiety treatment response found that conclusions depend on how response is defined and which analysis population is used \citep{Loerinc2015}.
A cross-condition synthesis found benefits that depended on the control condition, alongside limitations in trial quality and coverage \citep{Cuijpers2025}.
These findings motivate a critical review of nonresponse, comparative value and the scope of claims made for CBT.

The author's archived blog and Google Docs draft begin with a concern about reasoning through distress: cognitive restructuring may demand capacities that stress makes harder to use.
Stress-related prefrontal impairment provides a biological rationale for investigating this concern, but does not establish treatment failure \citep{Arnsten2009}.
The original drafts also ask whether exposure, behavioral practice and support explain benefits attributed to cognitive restructuring.
A depression component trial did not identify an advantage of the full cognitive-therapy package over its component treatments, providing a direct but bounded challenge to the necessity of the complete package \citep{Jacobson1996}.

The clinical evidence also constrains the critique.
Placebo-controlled anxiety trials report benefits of CBT, so a universal claim that CBT does not work is inconsistent with this retrieved evidence \citep{Carpenter2018}.
For depression, a large synthesis found benefit relative to controls but a small advantage over other psychotherapies that was not robust across most sensitivity analyses \citep{Cuijpers2023}.
These are different questions: treatment benefit, superiority over alternatives and support for a particular mechanism should not be collapsed into one verdict.

We therefore organize the review around distinct reasons for failure or overstatement across clinical conditions.
The scope includes incomplete response, comparator-sensitive effects, study bias, component attribution, executive demands, long-term functioning, acceptability and harms.
The current evidence map is a targeted starting synthesis; comprehensive cross-condition screening and quantitative extraction remain unfinished.
We preserve the author's skeptical argument and require each empirical claim to survive contradictory evidence and condition-specific assessment.

Adult attention-deficit/hyperactivity disorder (ADHD) is one case study within this broader review.
Existing reviews and trials describe benefits and limitations of CBT in adult ADHD, including executive-function and functioning outcomes \citep{Lopez2018,Safren2010,Liu2026,LopezPinar2026}.
A baseline predictor of outcome is distinct from a moderator of the randomized treatment contrast, and improvement in executive function is not itself evidence of baseline moderation \citep{Kraemer2002}.
Related older-adult studies of executive function and flexibility provide indirect evidence whose population and analyses must remain explicit \citep{Mohlman2005,Johnco2014}.

This manuscript contributes the amended protocol, a cited critical evidence map and a preliminary screen of the 21 records in the ADHD executive-function search slice.
The search baseline remains reproducible and is not relabeled as coverage of all CBT research.
Systematic selection and search reporting will follow the Preferred Reporting Items for Systematic Reviews and Meta-Analyses framework and its search extension \citep{Page2021,Rethlefsen2021}.
